\paragraph{Pucks.} We use the eight processed Slide-seqV2 pucks released by \citet{chen2021dissecting}: three wild-type (WT) mouse, three \emph{ob/ob} diabetic (DB) mouse, and two normal human (HU) pucks, each with a bead-by-gene count matrix, bead coordinates, and a per-bead NMFreg label over nine types (spermatogonia SPG, spermatocytes SPC, round spermatids RS, elongating/elongated spermatids ES, Sertoli, myoid, Leydig, endothelial, macrophage). We keep beads with a label and at least 100 UMIs, leaving 25,377 to 39,181 beads per puck (252,898 in total; Table~\ref{tab:homog}).

\paragraph{Cross-species gene space.} All pucks are restricted to one-to-one mouse--human orthologs from MGI \citep{baldarelli2024mouse} detected in every puck (13,573 genes), so a model trained on mouse can be applied to human without adaptation. Counts are library-size normalized and log-transformed \citep{wolf2018scanpy}. We select 2,000 highly variable genes on the WT pucks only \citep{stuart2019comprehensive}, then standardize each gene within each puck using that puck's own mean and standard deviation. This unsupervised per-puck standardization is the only batch handling we apply.

\paragraph{Task.} The primary task is four-class germ-cell stage annotation (SPG, SPC, RS, ES) on germ-cell beads, which are 81 to 90\% of beads. This is the layer the radial tubule axis defines, and the reproducible part of the NMFreg annotation: within-puck and cross-puck macro-F1 agree (0.71 vs.\ 0.70), whereas the rare somatic types (endothelial, macrophage, myoid) are near chance from a single bead's expression under any split. Six-class and nine-class variants are in the appendix.

\paragraph{Shift settings.} Training pucks are always disjoint from the test puck. \emph{WT$\to$WT}: train on two WT pucks, test on the third (three folds). \emph{WT$\to$ob/ob}: train on all three WT pucks, test on each DB puck. \emph{WT$\to$human}: train on all three WT pucks, test on each human puck; \emph{mouse$\to$human} adds the DB pucks to training. As in-distribution references for the shifted targets, \emph{ob/ob$\to$ob/ob} and \emph{human$\to$human} train on the other pucks of the same condition. The NMFreg labels were computed from expression alone, so a gain from spatial features cannot be an artefact of how the labels were made. For each puck we also compute neighbourhood homogeneity, the mean fraction of a bead's 15 nearest spatial neighbours that share its label; it is used only to describe the pucks, never as an input. Among germ-cell beads it is 0.53 to 0.56 on WT, 0.49 to 0.57 on \emph{ob/ob}, and 0.38 to 0.40 on human pucks.
